% Dot batch 87; source SHA-256: c619d13ca8fb077a894ab104684c849e26cdd2329765c40d7ef7b80a8ef13d17
% Archived source: dot_updates_20261124/batches/batch087/EGA第五卷_手写提纲与异扫同文.tex
% Original body lines: 14--46; integration map: integration/late_integration.json
\begingroup
\def\newscope{V-L-source}
\def\illis{〔字迹不清〕}
\def\gap{〔扫本缺损〕}
\head{来源补遗：邻接手稿与卷首异扫}
\egasection{法文扫描物理第150页：手写提纲}

\noindent 以下按扫描的物理页码定位。上半页是未署日期的条目草记，下半页另有倒写算式与划改。
刊载机构的目录将物理第148--150页归在致迪厄多内的书信名下；
但第150页既没有承接前页落款的句子，也没有与英译稿第112页对应的正文，
因此作为邻接手稿单列，不添入书信。\cmt{L18-087-01}{参见格罗滕迪克资料库\href{https://webusers.imj-prg.fr/~leila.schneps/grothendieckcircle/EGA/EGAVcontents.pdf}{第五卷预备笔记目录}。该目录提供档案分组，本页手写项目不能直接视为书信提出的同一版目录，也不是正式刊本章节目录。}

\subsec{正向条目}
\begin{description}
 \item[\S20] 关于 IV\illis{} 的补充；\illis{}；\illis{}，非双有理的\illis{}；以及爆破。\cmt{L18-087-02}{可辨首词为 \emph{complément}，后有形似罗马数字 IV 的标记，不能确定其后的编号。第二行可见 \emph{non} 及形如 \emph{birationnel} 的词根，末行可辨爆破一词；中间词和句法不强订。}
 \item[\S21] 除子，\illis{}，韦伊，皮卡尔；线性系，布劳尔--塞韦里。\cmt{L18-087-03}{在 \emph{Diviseurs} 后、\emph{Weil} 前有一组字迹可疑的缩记。此处不把它补成“卡蒂埃除子”，也不擅自指定“韦伊”“皮卡尔”两词的语法层级。}
 \item[\S22] 格拉斯曼簇（或概形）。\cmt{L18-087-04}{原文只写 \emph{Grassmanniennes}，未给出底概形或限定词。}
 \item[\S23] 超平面截面。
 \item[\S24] 无穷小扩张。\cmt{L18-087-05}{第23、24条分别为截面和无穷小扩张的缩记；第24条缩记不指定一个或多个扩张，本页无进一步定义。}
\end{description}

\subsec{旋转半周后阅读的笔记}
可辨重划的一行、一个另行划去的算式、一个向右箭头，以及下方若干短线连成的草图。
可辨题词“对偶”，其限定语和下接词句\illis{}。
未划去的主行可作下列残读；缺字不是数学运算符：
\[
 H^{\text{〔疑为 }n\text{〕}}(X)\ \simeq\ \sum
       \text{〔指标及算符不清〕}(X,D^{\text{〔疑为 }i\text{〕}}).
\]
左、下另有简写与连线，字迹不足以判定两端表达式或箭头类型。\cmt{L18-087-06}{保留残读只是为了定位手稿墨迹。级数的指标、上标、项的首符号和划去文字都未能核定，不能从此行推出任何上同调分解或对偶定理。包内另附旋转后的局部原图，供回查，不作为重排正文。}

\egasection{法文扫描第153页：卷首异扫}

这不是第四章第21节或书信的续页。它和法文扫描物理第1页
具有同一打字底文：下划线标题、16条拟定目录、目录后的编写说明及相同的半句截尾，
分行也基本对应。两页的边缘、反差和手写标记并不完全相同；不称为逐像素相同。
参考法文第1--2页及英译物理第2--3页可定位正文。
现行合订本已收相应目录与编写说明；对照译文见本版核查报告 L18-087，不计成新增一节。\cmt{L18-087-07}{第153页的原始编号栏多处不清。条号和部分词由法第1页及英第3页恢复；英第2页说明部分接续至法第2页。原档目录列出实存的第1--15节，草稿多列一条“改用线性系的语言”，二者不矛盾。}
\par\endgroup
